\documentclass{article}
\usepackage[utf8]{inputenc}
\usepackage{amsmath, amssymb}
\usepackage{siunitx}
\usepackage{booktabs}
\title{Detaljerte beregninger for EMP-generator (revidert)}
\author{}
\date{}

\begin{document}
\maketitle

\section{Innledning}
Dette dokumentet inneholder teoretiske beregninger for designet. Alle verdier er idealiserte og må verifiseres eksperimentelt. Vi skiller mellom matematisk korrekte utledninger, antatte parametere og estimater med usikkerhet.

\section{Hovedkondensator og spole}

\subsection{Spesifikasjoner}
\begin{itemize}
\item Kondensator: $C = 1000\,\mu\text{F}$, merkespenning $450\,\text{V}$ (lades til $V_0 = 400\,\text{V}$ for sikkerhetsmargin).
\item Spole: Flat spiral, $N = 6$ vindinger, indre diameter $20\,\text{mm}$, ytre diameter $100\,\text{mm}$, tråd: $10\,\text{AWG}$ ($2.588\,\text{mm}$ diameter).
\end{itemize}

\subsection{Induktans for flat spiralspole}
Wheelers formel for en flat spiralspole i luft, med dimensjoner i tommer og resultat i µH, er:
\begin{equation}
L = \frac{r^2 N^2}{8r + 11w},
\end{equation}
der $r$ er middelradius i tommer, $w$ er viklingens bredde i tommer.

Konvertering: $r = 30\,\text{mm} = 1.181\,\text{tommer}$, $w = 40\,\text{mm} = 1.575\,\text{tommer}$.
\begin{equation}
L = \frac{(1.181)^2 \cdot 6^2}{8 \cdot 1.181 + 11 \cdot 1.575} = \frac{1.395 \cdot 36}{9.448 + 17.325} = \frac{50.22}{26.773} \approx 1.876\,\mu\text{H}.
\end{equation}
Vi bruker $L = 1.9\,\mu\text{H}$ som et teoretisk estimat. Målt verdi kan avvike.

\subsection{Energi lagret i kondensator}
\begin{equation}
E = \frac{1}{2} C V_0^2 = 0.5 \cdot 1000 \times 10^{-6} \cdot (400)^2 = 80\,\text{J}.
\end{equation}

\section{Toppstrøm og pulsform}

\subsection{Karakteristisk impedans}
\begin{equation}
Z_0 = \sqrt{\frac{L}{C}} = \sqrt{\frac{1.9 \times 10^{-6}}{1000 \times 10^{-6}}} = \sqrt{1.9 \times 
